\documentclass[11pt]{article}
\usepackage[a4paper,margin=25mm]{geometry}
\usepackage{amsmath,amssymb,amsthm,hyperref}
\newtheorem{theorem}{Theorem}\newtheorem{lemma}{Lemma}
\newcommand{\E}{\mathbb E}
\title{The square-root interior endpoint profile}
\author{Complete theorem; independently reviewed}
\date{30 September 2026}
\begin{document}\maketitle

\begin{theorem}\label{main}
Suppose $p_1,\ldots,p_m$ are nondecreasing $[0,1]$-valued columns
on a finite ordered set with positive probability row weights $w$,
and
\[
 \E_w\prod_{j\in S}p_j=\frac1{|S|+1}\qquad(S\subseteq[m]).
\]
Put $u_i=1-p_i$ and $s=\sum_i u_i$. For every fixed
$0<\delta<L<\infty$, uniformly over rows with $\delta\le s\le L$,
\[
 \max_i|m u_i-s|\le C_{\delta,L}m^{-1/2}.
\]
\end{theorem}

We use the exact marked identity and notation from
\path{marked_cancellation_identity.tex}, and the common-coordinate
bookkeeping from \path{sharp_interior_endpoint_profile.tex}. The
present argument replaces its infinite Laguerre kernel by a positive
finite polynomial. We give all changes to the measure and maximum
arguments explicitly; Gaussian tails are not assumed here.

The new analytic inputs are
\path{../rational_designs/asymptotic/finite_jackson_endpoint_kernel.tex}
(including its global low-order off-annulus estimate) and
\path{../bernstein_walsh_polarization_tail.tex}.
As a prior estimate we use any independently proved bound
\begin{equation}\label{prior}
 \max_i|m u_i-s|\le R_m\le m^{-\alpha},\qquad\alpha>0,
\end{equation}
on each fixed bounded $s$-interval. The earlier polynomial-rate
profile theorem suffices; the stronger logarithmic square-root
theorem is not needed as an input.

\paragraph{Parameters and local comparisons.}
Fix a large error exponent $A>10$. Fix the positive integer $r$ in
the Jackson kernel sufficiently large that
\begin{equation}\label{rchoice}
 2r\alpha>A+10,\qquad r>A+10.
\end{equation}
The polynomial exponents absorbed here are fixed independently of
$r$. Constants, as distinct from powers of $m$, may depend on $r$.
Next choose the degree ratio sufficiently small for both the global
Bernstein--Walsh theorem and the Jackson off-annulus geometric sum
(the latter constants depend on $r$), and finally choose its differential cutoff
$J=B\log m$. For the Jackson kernel $Q_m(x,a)$ let
\[
 b=\sqrt\ell/\nu\asymp m^{-1/2},\qquad
 E_b(x,a)=b^{-1}(1+|x-a|/b)^{-2r}.
\]
Here $\ell=\lfloor(m-1)/4\rfloor$, and the integer $\nu\asymp m$
has fixed sufficiently small ratio, after $r$ is fixed.

For marked column $i$, put $M=m-1$, $N=M-\ell$, $z_j=\ell u_j$,
and use
\[
 x_i=M^{-1}\sum_{j\ne i}z_j,\qquad h_i=z_i-x_i,\qquad
 d\mu_i=(\ell+1)\E_B\prod_{j\in B}p_j\,dw,
\]
where $B$ ranges uniformly over the $\ell$-subsets of the other
$M$ indices. Each $\mu_i$ is a probability measure. The reference
measure has density
\[
 d\beta_\ell(x)=\frac{\ell+1}{\ell}(1-x/\ell)^\ell
                  \boldsymbol1_{[0,\ell]}(x)\,dx.
\]
Let $H=\max_jz_j-\min_jz_j$.

The exact polynomial identities, global high-order tail theorem,
Jackson off-annulus estimate, and compact subset comparisons give,
for centers in any fixed positive compact interval,
\begin{align}
 |\mu_i(Q_m(x_i,a))-\beta_\ell(Q_m(\cdot,a))|
 &\le \frac C{mb^2}\int_{x_i\in W}H^2E_b(x_i,a)\,d\mu_i+m^{-A},
 \label{unmarked}\\
 |\mu_i(h_iQ_m(x_i,a))|
 &\le\frac C m\int_{x_i\in W}
 H^2(1+b^{-1}+Hb^{-2})E_b(x_i,a)\,d\mu_i+m^{-A}.
 \label{marked}
\end{align}
Here $W$ is a sufficiently large fixed positive row annulus, and
constants or the envelope scale can change by fixed factors.
To see the order of these operations, first apply the exact marked
and unmarked identities to the polynomial $Q_m$, whose degree is
at most a small fixed multiple of $m$. Remove differential orders
above $J$ by the global bounded-polynomial tail theorem. At orders
up to $J$, discard rows outside $W$ using the Jackson angular
localization theorem. On $W$ use the finite centered differential
sum and the success-weighted subset Taylor comparison. The prior
estimate makes $H$ small there, and finite balance ensures every
marked value is bounded whenever $x_i$ is bounded.
The scalar terms outside $W$ are also harmless: the global scalar
kernel tail is $O(b^{2r-1})$, each row factor is polynomially bounded,
and \eqref{rchoice} absorbs their fixed powers of $m$.

\paragraph{Mass absorption with polynomial tails.}
We record why the earlier Gaussian mass argument still applies.
On $W$, \eqref{prior} gives $H\le CR_m$, so \eqref{unmarked}
has right side at most
\[
 \epsilon_m\mu_i(E_b(x_i,a))+m^{-A},\qquad
 \epsilon_m=CR_m^2/(mb^2)\longrightarrow0.
\]
The Jackson kernel is at least $c/b$ on $|x-a|\le c_r b$, and its
reference integral is bounded above and below by positive constants.
Fix an outer positive annulus $I$ and let
$\phi(a)=\operatorname{dist}(a,\partial I)^2$.
Maximize $\phi(a)\mu_i([a-c_r b,a+c_r b])/b$ over $a\in I$.
If this supremum is at least one, the total-mass bound implies that
an approximate maximizer has boundary distance $d\ge c\sqrt b$.
On its $d/2$ neighborhood the weights $\phi$ are comparable, and
summing the envelope over intervals of length $c_r b$ bounds the
near part by a constant times the supremum divided by $\phi(a)$.
The far part is at most
\[
 Cb^{-1}(d/b)^{-2r}=O(b^{r-1}),
\]
which tends to zero. Thus the same weighted-supremum absorption
proves a uniform local interval bound. On any strictly smaller
annulus, it yields
\begin{equation}\label{mass}
 \mu_i(|x_i-a|\le b)\le Cb,\quad
 c\le\mu_i(Q_m(x_i,a))\le C,\quad \mu_i(E_b(x_i,a))\le C.
\end{equation}
Summation over length-$b$ intervals also gives, for $T\ge1$,
\begin{equation}\label{polytail}
 \int_{|x_i-a|\ge Tb}E_b(x_i,a)\,d\mu_i
       \le C T^{1-2r}+O(m^{-A}),
\end{equation}
as long as the local interval argument is applied in a fixed larger
annulus. Tails beyond that larger annulus are bounded directly by
the total mass and the global scalar kernel bound. By choosing
$r$ as in \eqref{rchoice}, those tails have any required fixed
inverse-power accuracy in this proof.

\paragraph{The weighted maximum and its far tail.}
Use the common coordinate and deviations
\[
 y=\ell s/m,\quad h_i^*=z_i-y,\quad
 x_i=y-h_i^*/M,\quad h_i=(m/M)h_i^*.
\]
Choose a fixed positive annulus $I=[a_-,a_+]$ whose strictly smaller
interior contains the desired $y$-range. The mass bounds are first
proved on an even larger annulus. Let
\[
 D=\max_{i,\,y\in I}\operatorname{dist}(y,\partial I)|h_i^*(y)|.
\]
If $D>C_*b$, select a maximizing row $y_0$ and column $i$, and write
$\Delta=|h_i^*(y_0)|$, $d_0=\operatorname{dist}(y_0,\partial I)$.
Then
\[
 d_0/b=D/(b\Delta)\ge C_*/R_m\ge C_*m^\alpha,
 \qquad \Delta\ge c_I C_*b.
\]
In the $d_0/2$ row neighborhood, weighted maximality gives
$|h_j^*|\le C\Delta$ for every $j$, hence $H\le C\Delta$.
The deleted coordinate differs from $y$ by $O(R_m/m)=o(b)$.
Any row entering the fixed deleted-coordinate neighborhood has
bounded full $s$ by finite balance, so this estimate applies before
restricting to the common-coordinate neighborhood.

First discard rows outside a fixed positive enlargement by the global
Jackson off-annulus and scalar tail estimates already used in
\eqref{unmarked}--\eqref{marked}. The scalar Jackson tail there is
$O(b^{2r-1}(1+x)^{-r})$; we do not bound it by $E_b$ globally.
Within that fixed enlargement, $Q_m\le CE_b$, and every far
contribution beyond a fixed multiple of $d_0$ has envelope
at most
\[
 C b^{-1}(d_0/b)^{-2r}\le C m^{1/2-2r\alpha}.
\]
The measures have mass one and all scalar prefactors used here are
bounded by $Cm^4$ (in fact much smaller bounds suffice).
Consequently \eqref{rchoice} makes all those contributions smaller
than $m^{-A}$. This bound has an exponent independent of $r$ except
for the favorable decay term $-2r\alpha$.

\paragraph{The sign test.}
If the selected deviation is positive, set
$a=x_i(y_0)+A_0b$; for a negative deviation use the minus sign.
Choose fixed $B_0$ large, then $A_0=2B_0$.
On $|x_i-a|\le B_0b$ the row lies strictly on the favorable side
of the selected one. Monotonicity therefore gives, in the positive
case,
\[
 h_i(x_i)\ge h_i(x_i(y_0))-(x_i-x_i(y_0))\ge\Delta/2,
\]
provided $C_*$ is sufficiently large. The negative case is identical
with signs reversed. Strict sidedness handles flat rows.
By \eqref{mass}--\eqref{polytail}, this window captures a positive
amount of $Q_m$-mass, while the near opposite-sign contribution is
at most $C\Delta B_0^{1-2r}$. Choose $B_0$ large enough to absorb it.
The farther contribution is negligible by the preceding paragraph.
Thus
\[
 |\mu_i(h_iQ_m(x_i,a))|\ge c\Delta.
\]
But \eqref{marked}, the local bound $H\le C\Delta$, and
\eqref{mass} give
\[
 |\mu_i(h_iQ_m(x_i,a))|
 \le C(\Delta^2/\sqrt m+\Delta^3)+m^{-A}=o(\Delta),
\]
because $\Delta\le R_m\to0$ and $\Delta\ge c_I C_*m^{-1/2}$.
This contradiction proves $D=O(b)$. The cutoff is bounded below
on the target range, so $h_i^*=O(m^{-1/2})$ there.
Since $m u_i-s=(m/\ell)h_i^*$, Theorem~\ref{main} follows.
\end{document}
